\chapter{Constructing a Balanced, Rule-Mapped Pilot Sample (Day 2)}
\label{chap:day02}

\section{Motivation}

With the environment established in Chapter~1, the pilot needed a fixed, finite sample to evaluate
against, rather than an open-ended stream of dataset rows. ConstructionSite 10k's test split
contains thousands of images covering four safety conditions of interest --- PPE violations, fall
hazards, struck-by risk, and full compliance --- in proportions set by the dataset itself, not by
the pilot's evaluation needs. Drawing a plain random sample from that split risks a sample
dominated by whichever class happens to be most common, leaving the rarer classes too thin to
support any later per-class analysis. Day 2's task was therefore to construct a sample that is
(i) balanced across the four classes the pilot cares about, (ii) explicitly mapped to the specific
safety rule each sample is meant to test, and (iii) reproducible from a fixed seed, before any
inference or evaluation work began.

\section{Implementation}

Three functions, added to \texttt{data.py} in commit \texttt{c76a642}, carry Day 2's logic, together
with a new \texttt{RULE\_QUERIES} table in \texttt{prompts.py} and a driver script,
\texttt{scripts/02\_select\_samples.py}, that writes the resulting sample to
\texttt{pilot\_samples.csv}.

\texttt{classify\_image(row)} maps a dataset row's rule-violation fields to one of four classes.
A row may violate zero, one, or several of the four safety rules; \texttt{classify\_image} returns
a single \texttt{primary\_class}, taken as the highest-priority violated class according to
\texttt{config.py}'s \texttt{CLASS\_PRIORITY = ["ppe\_violation", "fall\_hazard",
"struck\_by\_risk", "compliant"]}. A row that violates both rule~1 and rule~3, for example, is
classified \texttt{ppe\_violation} rather than \texttt{fall\_hazard}, because \texttt{ppe\_violation}
sits earlier in the priority list. Separately, \texttt{violated\_rule\_ids} records every rule the
row actually violates, not only the one that determined \texttt{primary\_class}.

\texttt{assign\_rule\_id} then picks the one rule each sample is tested against in Days 3--10.
\texttt{config.py}'s \texttt{RULE\_TO\_CLASS} maps \texttt{rule\_1\_violation} to
\texttt{ppe\_violation}, \texttt{rule\_2\_violation} and \texttt{rule\_3\_violation} both to
\texttt{fall\_hazard}, and \texttt{rule\_4\_violation} to \texttt{struck\_by\_risk}. For a violation
sample, the assigned rule is the one unambiguous violated rule that maps to its
\texttt{primary\_class}. Compliant samples have no ground-truth rule to assign, since by
definition they violate none; these are instead round-robined evenly across
\texttt{COMPLIANT\_ROUND\_ROBIN = ["rule\_1", "rule\_2", "rule\_3", "rule\_4"]}. This round-robin is
a deliberate design choice: it gives balanced "model correctly says compliant" coverage for each
rule, rather than testing the same rule against every compliant image and leaving the other three
rules with no compliant counterexamples at all.

\texttt{select\_balanced\_sample} performs the actual selection in a single streaming pass over the
dataset. For each of the four classes, a reservoir of candidate samples is kept, capped at
\texttt{n\_per\_class * 10} (500 candidates) to bound memory while streaming. Once reservoirs are
populated, the final sample is a seeded (\texttt{SEED = 42}) random draw of \texttt{n\_per\_class}
(\texttt{SAMPLES\_PER\_CLASS = 50}) from each class's reservoir, making the selection reproducible
from the same seed and dataset stream.

\texttt{prompts.py} additionally introduces \texttt{RULE\_QUERIES}, two phrasings per presence-based
rule used later to query Florence-2: \texttt{rule\_1: ["hard hat", "high-visibility vest"]},
\texttt{rule\_2: ["safety harness", "fall-protection lanyard"]}, and \texttt{rule\_3: ["guardrail",
"edge protection barrier"]}. \texttt{PERSON\_PHRASE = "worker"} and
\texttt{RULE\_4\_PROXIMITY\_PAIR = ("worker", "excavator")} support rule~4's proximity-based check,
which is structurally different from the presence checks used for rules 1--3. \texttt{RULE\_QUERIES}
is not yet consumed on Day 2; it is written here so that Day 3's inference step has a single,
already-defined source of query phrasings to draw from.

\section{Results}

\texttt{scripts/02\_select\_samples.py} streamed the ConstructionSite test split once and wrote the
selected sample to \texttt{pilot\_samples.csv}. The result is 163 total rows, not the 200 originally
planned (50 per class across four classes):

\begin{center}
\begin{tabular}{lr}
\toprule
\textbf{Class} & \textbf{Count} \\
\midrule
\texttt{ppe\_violation} & 50 \\
\texttt{fall\_hazard} & 50 \\
\texttt{compliant} & 50 \\
\texttt{struck\_by\_risk} & 13 \\
\midrule
\textbf{Total} & 163 \\
\bottomrule
\end{tabular}
\end{center}

Three of the four classes filled to the planned target of 50. \texttt{struck\_by\_risk} filled to
only 13, capped by genuine scarcity in the underlying dataset rather than by any defect in the
sampling procedure: the entire 3{,}004-image test split contains only 13 \texttt{struck\_by\_risk}
examples in total, a fact confirmed directly by an exhaustive scan of the full split rather than
assumed from the reservoir's behavior.

\section{Problems Encountered and Resolutions}

The shortfall in \texttt{struck\_by\_risk} samples was the one substantive problem encountered on
Day 2. It surfaced as \texttt{select\_balanced\_sample}'s reservoir for that class filling to only
13 candidates regardless of how much of the stream was consumed. Rather than treat this as a
sampling bug to be patched, it was investigated directly: a full scan of the test split confirmed
that 13 is the true population count for \texttt{struck\_by\_risk} in that split, not an artifact
of the streaming or reservoir logic. \emph{Resolution}: proceed with \texttt{n=13} for
\texttt{struck\_by\_risk} rather than block the pilot on a class the dataset cannot supply at the
planned size. This shortfall is flagged explicitly as a limitation here, and it recurs in every
later chapter's \texttt{struck\_by\_risk} numbers.

\section{Standard vs. Non-Standard Choices}

The sampling strategy is, in its broad outline, standard practice: stratified, class-balanced
sampling implemented via per-class reservoir sampling with a fixed seed, so the same sample can be
reconstructed from the same dataset stream. This part of the design follows well-established
sampling methodology and is not particular to this pilot.

The round-robin assignment of rules to compliant samples, by contrast, is non-standard: it is a
domain-specific balancing choice, motivated by the absence of a natural ground-truth rule for
compliant images, and is not a pattern found in generic sampling libraries. It was introduced here
specifically so that each of the four safety rules receives roughly equal compliant-sample coverage
for later "correctly identifies compliance" evaluation, rather than leaving that coverage to chance.

\section{Implications for the Paper}

The 13-vs-50 class imbalance in \texttt{struck\_by\_risk} is a property of the ConstructionSite 10k
dataset itself, not a limitation of the sampling method used to draw from it, and should be reported
as such in any paper's Dataset or Limitations section. Because the imbalance was confirmed by an
exhaustive scan rather than inferred from the pilot's own reservoir behavior, it can be stated with
confidence rather than hedged. This scarcity directly motivates an open question raised again in
Chapter~14: whether a future iteration of this work should scale up data collection specifically for
underrepresented classes such as \texttt{struck\_by\_risk}, or instead accept and explicitly correct
for the imbalance during evaluation. With a fixed, balanced, rule-mapped sample of 163 images now in
hand, the pilot's next step is to run baseline inference over that sample.
